% SEGMENT OLP-0605-S01
% Part: methods
% Chapter: proofs
% Section: using-definitions

\documentclass[../../../include/open-logic-section]{subfiles}

\begin{document}

% SOURCE CORRECTION TA-MTH-001: பகுதியின் mth அடையாளத்துக்கு மாறாக இருந்த mod ஐச் சரிசெய்தோம்.
\olfileid{mth}{prf}{def}

\olsection{வரையறைகளைப் பயன்படுத்துதல்}

நிறுவலில் பயன்படக்கூடிய எல்லா வரையறைகளையும் அறிந்திருக்க
வேண்டும் என்றும், அவற்றைச் சரியாகப் பயன்படுத்த வேண்டும்
என்றும் முன்பு கூறினோம். இது மிகவும் முக்கியமானது; சற்று
விரிவாகப் பார்ப்போம். பண்புகளையும் தொடர்புகளையும் சுருக்கமாகக்
கூற வரையறைகள் உதவுகின்றன. அறிமுகப்படுத்தப்படும் சுருக்கப் பெயர்
\emph{வரையறுக்கப்படுவது}; அது சுருக்கிக் குறிக்கும் விளக்கம்
\emph{வரையறுக்கும் பகுதி}. நிறுவலில் வரையறுக்கப்பட்ட சொல்லை
எப்படி அறிமுகப்படுத்தினோம் என்பதற்கு அடிக்கடி திரும்ப வேண்டும்.
ஏனெனில், நிறுவலை முன்னெடுக்க வரையறுக்கும் பகுதியின், அதாவது
வரையறுக்கப்பட்ட சொல்லின் விரிவான வடிவத்தின், தருக்க அமைப்பைப்
பயன்படுத்த வேண்டியிருக்கும். வரையறைகளை விரித்துரைப்பதன்
மூலம் தருக்கச் செயற்பாடு நடக்கும் மையத்தை அடைகிறீர்கள்.

% SEGMENT OLP-0605-S02
ஓர் எடுத்துக்காட்டுடன் தொடங்குவோம். பின்வருவதை நிறுவ
விரும்புகிறீர்கள் என்க:

\begin{prop}
எந்தக் கணங்கள் $A$, $B$ க்கும் $A \cup B = B \cup A$.
\end{prop}

நிறுவலைத் தொடங்குவதற்கே இரண்டு கணங்கள் ஒன்றாக இருப்பது
என்னவென்று அறிய வேண்டும்; அதாவது, அந்தச் சமன்பாட்டில்
கணங்களுக்கான ``$=$'' குறி என்ன பொருள் தருகிறது என்று அறிய
வேண்டும். கணங்கள் ஒரே !!{element}s ஐக் கொண்டிருக்கும்போது
அவை ஒன்றாக இருப்பதாக வரையறுக்கப்படுகின்றன. எனவே விரிக்க
வேண்டிய வரையறை இதுதான்:

\begin{defn}
கணங்கள் $A$, $B$ ஆகியவை \emph{ஒன்றானவை}, $A = B$, எனில்,
எனில் மட்டுமே, $A$ இன் ஒவ்வோர் !!{element} உம் $B$ இன்
!!a{element} ஆகவும், மறுதிசையிலும் இருக்கும்.
\end{defn}

இந்த வரையறையில் $A$, $B$ ஆகியவை எந்தக் கணத்துக்கும்
பொருந்தும் இடக்குறிகளாகப் பயன்படுத்தப்படுகின்றன. இதன்
\emph{வரையறுக்கப்படுவது} ``$A = B$'' என்ற கோவை;
$A = B$ மெய்யாக இருப்பதற்கான நிபந்தனையை இது தருகிறது.
``$A$ இன் ஒவ்வோர் !!{element} உம் $B$ இன்
!!a{element} ஆகவும், மறுதிசையிலும் இருக்கும்'' என்ற
நிபந்தனையே \emph{வரையறுக்கும் பகுதி}.\footnote{இந்தக்
  குறிப்பிட்ட வழக்கில், சற்றுக் குழப்பமூட்டும் விதமாக,
  $A = B$ எனில் $A$, $B$ ஆகியவை உண்மையில் ஒரே கணம்;
  இடப்புறமும் வலப்புறமும் அதற்கு வெவ்வேறு எழுத்துகளைப்
  பயன்படுத்துகிறோம். ஆனால் அக்கணத்தை அடையாளப்படுத்தும்
  வழிகள் வேறுபடலாம்; அதனால் இந்த வரையறை வெறுமையானதல்ல.}
இந்த நிபந்தனை நிறைவேறினால், நிறைவேறினால் மட்டுமே, $A = B$
மெய் என்று வரையறை கூறுகிறது. இதைத்தான் ``எனில், எனில்
மட்டுமே'' என்று சுருக்கமாகச் சொல்கிறோம்.

% SEGMENT OLP-0605-S03
வரையறையைப் பயன்படுத்தும்போது, அதிலுள்ள $A$, $B$
ஆகியவற்றை உங்கள் குறிப்பிட்ட வழக்கின் பொருள்களுடன்
பொருத்த வேண்டும். இங்கு $A \cup B = B \cup A$ மெய்யாக
இருக்க வேண்டுமானால், ஒவ்வொரு $z \in A \cup B$ உம்
$B \cup A$ இல் இருக்க வேண்டும்; மறுதிசையிலும் அதேபோல்
இருக்க வேண்டும். முன்மொழிவிலுள்ள $A \cup B$ வரையறையின்
$A$ ஆகவும், $B \cup A$ அதன் $B$ ஆகவும் அமைகின்றன.
வரையறையிலும் நிறுவப்படும் முன்மொழிவிலும் $A$, $B$ என்ற
எழுத்துகளே வெவ்வேறு பணிகளில் வருவதால், அவற்றைக்
கலக்காமல் கவனமாக இருக்க வேண்டும். எடுத்துக்காட்டாக,
$A$ இன் ஒவ்வோர் !!{element} உம் $B$ இன்
!!a{element} ஆகவும், மறுதிசையிலும் உள்ளது என்பதைக்
காட்டினால் முன்மொழிவு நிறுவப்பட்டுவிடும் என்று கருதுவது
தவறு. அது $A = B$ என்பதைத்தான் காட்டும்;
$A \cup B = B \cup A$ என்பதைக் காட்டாது. மேலும்,
$A$, $B$ எந்த இரு கணங்களாகவும் இருக்கலாம். அவற்றைப்
பற்றி எதுவும் எடுக்கப்படாதபோது $A$, $B$ வெவ்வேறாகவே
இருக்கலாம்; ஆகவே அப்படி முயன்று வெகுதூரம் செல்ல முடியாது.

% SEGMENT OLP-0605-S04
நிறுவலில் ஒன்றிப்பு போன்ற கணக் கோட்பாட்டுக் கருத்துகளைக்
கையாள்கிறோம். எனவே நிறுவல் எப்படி முன்னேற வேண்டும் என்று
புரிந்துகொள்ள $\cup$ குறியின் பொருளையும் அறிந்திருக்க வேண்டும்.
ஒரு வரையறையை விரிக்கும்போது மேலும் சில வரையறைகளையும்
விரிக்க வேண்டியிருக்கலாம். எடுத்துக்காட்டாக, $A \cup B$
என்பது $\Setabs{z}{z \in A \text{ அல்லது } z \in B}$
என்று வரையறுக்கப்படுகிறது. ஆகவே $x \in A \cup B$
என்பதை நிறுவ விரும்பினால், $\cup$ இன் வரையறையை
விரிப்பதால் $x \in \Setabs{z}{z \in A \text{ அல்லது } z \in B}$
என்பதைக் காட்ட வேண்டும் என்று தெரிகிறது. மேலும்,
$x \in \Setabs{z}{\dots z\dots}$ எனில், எனில் மட்டுமே,
$\dots x\dots$ என்பதையும் நினைவில் கொள்ள வேண்டும்.
எனவே $\Setabs{z}{\dots z \dots}$ குறியீட்டின் வரையறையையும்
விரித்தால், காட்ட வேண்டியது $x \in A$ அல்லது $x \in B$
என்பதே. அதாவது, ``$A \cup B$ இன் ஒவ்வோர் !!{element} உம்
$B \cup A$ இன் !!a{element} ஆகும்'' என்பதன் முழுப் பொருள்:
``ஒவ்வொரு $x$ க்கும், $x \in A$ அல்லது $x \in B$ எனில்,
$x \in B$ அல்லது $x \in A$.'' முன்மொழிவின்
வரையறைகளை முற்றிலும் விரித்தால், நாம் காட்ட வேண்டியது
பின்வருவதாகிறது:

\begin{prop}
எந்தக் கணங்கள் $A$, $B$ க்கும்: (a) ஒவ்வொரு $x$ க்கும்,
$x \in A$ அல்லது $x \in B$ எனில், $x \in B$ அல்லது
$x \in A$; (b) ஒவ்வொரு $x$ க்கும், $x \in B$ அல்லது
$x \in A$ எனில், $x \in A$ அல்லது $x \in B$.
\end{prop}

% SEGMENT OLP-0605-S05
வரையறைகளை விரிப்பது நிறுவலை உருவாக்குவதற்கு இன்றியமையாதது.
அதைச் சரியாகச் செய்வது சில நேரம் கடினம்: வரையறையிலுள்ள
மாறிகளையும் நிறுவப்படும் கூற்றிலுள்ள சொற்களையும் வேறுபடுத்தி,
சரியாகப் பொருத்த வேண்டும். வெற்றிபெற, கேள்வி என்ன கேட்கிறது
என்றும் அதில் பயன்படுத்தப்படும் எல்லாச் சொற்களின் பொருளும்
என்னவென்றும் அறிய வேண்டும். ஒன்றுக்கு மேற்பட்ட வரையறைகளை
அடிக்கடி விரிக்க வேண்டியிருக்கும். கீழே உள்ளவை போன்ற எளிய
நிறுவல்களில், வரையறைகளிலிருந்தே விடை கிட்டத்தட்ட உடனடியாகப்
பின்தொடரும். எல்லா நேரமும் இவ்வளவு எளிதாக இருக்காது.

\begin{prob}
$A \cap B \neq \emptyset$ என்பதை நிறுவுமாறு கேட்கப்படுகிறது
என்க. இங்கு வரும் எல்லா வரையறைகளையும் விரிக்கவும்;
அதாவது, ``$\cap$'', ``='', ``$\emptyset$'' ஆகியவற்றைக்
குறிப்பிடாமல் இதே கருத்தை மீண்டும் கூறவும்.
\end{prob}

\end{document}
